%=========================================================================
%  Séance 1 : Environnement et anatomie des données
%  Compilation :  pdflatex 01-intro.tex   (deux fois)
%=========================================================================
\documentclass[10pt,aspectratio=169]{beamer}
%=========================================================================
%  Préambule commun — Analyse des données, L3 Économie
%  Appelé par chaque séance via %=========================================================================
%  Préambule commun — Analyse des données, L3 Économie
%  Appelé par chaque séance via %=========================================================================
%  Préambule commun — Analyse des données, L3 Économie
%  Appelé par chaque séance via \input{preambule-cy}
%=========================================================================
\usetheme[progressbar=frametitle,numbering=fraction,block=fill]{metropolis}

\usepackage[T1]{fontenc}
\usepackage[utf8]{inputenc}
\usepackage[french]{babel}
%\usepackage{lmodern}  % absent de cet environnement de compilation ; a decommenter
\usepackage{booktabs}
\usepackage{listings}
\usepackage{xcolor}

%-------------------------------------------------------------------------
%  PALETTE — remplacer par les valeurs de votre charte CY habituelle
%-------------------------------------------------------------------------
\definecolor{CYprincipal}{HTML}{1B2A4A}   % couleur dominante (titres, barres)
\definecolor{CYaccent}{HTML}{C8102E}      % couleur d'accentuation
\definecolor{CYgris}{HTML}{55637C}        % texte secondaire
\definecolor{CYclair}{HTML}{EEF1F6}       % fonds de blocs

\setbeamercolor{palette primary}{bg=CYprincipal,fg=white}
\setbeamercolor{progress bar}{fg=CYaccent}
\setbeamercolor{frametitle}{bg=CYprincipal,fg=white}
\setbeamercolor{alerted text}{fg=CYaccent}
\setbeamercolor{block title}{bg=CYclair,fg=CYprincipal}
\setbeamercolor{block body}{bg=CYclair!60,fg=black}

%-------------------------------------------------------------------------
%  Mise en forme du code Python
%-------------------------------------------------------------------------
\lstdefinestyle{py}{
  language=Python,
  basicstyle=\ttfamily\small,
  keywordstyle=\color{CYaccent},
  stringstyle=\color{CYprincipal},
  commentstyle=\color{CYgris}\itshape,
  showstringspaces=false,
  breaklines=true,
  frame=none,
  literate={é}{{\'e}}1 {è}{{\`e}}1 {à}{{\`a}}1 {ç}{{\c c}}1 {ê}{{\^e}}1 {ô}{{\^o}}1 {û}{{\^u}}1 {î}{{\^i}}1
}
\lstset{style=py}

\author{Stefania Marcassa}
\institute{CY Cergy Paris Université --- L3 Économie}
\date{}

%=========================================================================
\usetheme[progressbar=frametitle,numbering=fraction,block=fill]{metropolis}

\usepackage[T1]{fontenc}
\usepackage[utf8]{inputenc}
\usepackage[french]{babel}
%\usepackage{lmodern}  % absent de cet environnement de compilation ; a decommenter
\usepackage{booktabs}
\usepackage{listings}
\usepackage{xcolor}

%-------------------------------------------------------------------------
%  PALETTE — remplacer par les valeurs de votre charte CY habituelle
%-------------------------------------------------------------------------
\definecolor{CYprincipal}{HTML}{1B2A4A}   % couleur dominante (titres, barres)
\definecolor{CYaccent}{HTML}{C8102E}      % couleur d'accentuation
\definecolor{CYgris}{HTML}{55637C}        % texte secondaire
\definecolor{CYclair}{HTML}{EEF1F6}       % fonds de blocs

\setbeamercolor{palette primary}{bg=CYprincipal,fg=white}
\setbeamercolor{progress bar}{fg=CYaccent}
\setbeamercolor{frametitle}{bg=CYprincipal,fg=white}
\setbeamercolor{alerted text}{fg=CYaccent}
\setbeamercolor{block title}{bg=CYclair,fg=CYprincipal}
\setbeamercolor{block body}{bg=CYclair!60,fg=black}

%-------------------------------------------------------------------------
%  Mise en forme du code Python
%-------------------------------------------------------------------------
\lstdefinestyle{py}{
  language=Python,
  basicstyle=\ttfamily\small,
  keywordstyle=\color{CYaccent},
  stringstyle=\color{CYprincipal},
  commentstyle=\color{CYgris}\itshape,
  showstringspaces=false,
  breaklines=true,
  frame=none,
  literate={é}{{\'e}}1 {è}{{\`e}}1 {à}{{\`a}}1 {ç}{{\c c}}1 {ê}{{\^e}}1 {ô}{{\^o}}1 {û}{{\^u}}1 {î}{{\^i}}1
}
\lstset{style=py}

\author{Stefania Marcassa}
\institute{CY Cergy Paris Université --- L3 Économie}
\date{}

%=========================================================================
\usetheme[progressbar=frametitle,numbering=fraction,block=fill]{metropolis}

\usepackage[T1]{fontenc}
\usepackage[utf8]{inputenc}
\usepackage[french]{babel}
%\usepackage{lmodern}  % absent de cet environnement de compilation ; a decommenter
\usepackage{booktabs}
\usepackage{listings}
\usepackage{xcolor}

%-------------------------------------------------------------------------
%  PALETTE — remplacer par les valeurs de votre charte CY habituelle
%-------------------------------------------------------------------------
\definecolor{CYprincipal}{HTML}{1B2A4A}   % couleur dominante (titres, barres)
\definecolor{CYaccent}{HTML}{C8102E}      % couleur d'accentuation
\definecolor{CYgris}{HTML}{55637C}        % texte secondaire
\definecolor{CYclair}{HTML}{EEF1F6}       % fonds de blocs

\setbeamercolor{palette primary}{bg=CYprincipal,fg=white}
\setbeamercolor{progress bar}{fg=CYaccent}
\setbeamercolor{frametitle}{bg=CYprincipal,fg=white}
\setbeamercolor{alerted text}{fg=CYaccent}
\setbeamercolor{block title}{bg=CYclair,fg=CYprincipal}
\setbeamercolor{block body}{bg=CYclair!60,fg=black}

%-------------------------------------------------------------------------
%  Mise en forme du code Python
%-------------------------------------------------------------------------
\lstdefinestyle{py}{
  language=Python,
  basicstyle=\ttfamily\small,
  keywordstyle=\color{CYaccent},
  stringstyle=\color{CYprincipal},
  commentstyle=\color{CYgris}\itshape,
  showstringspaces=false,
  breaklines=true,
  frame=none,
  literate={é}{{\'e}}1 {è}{{\`e}}1 {à}{{\`a}}1 {ç}{{\c c}}1 {ê}{{\^e}}1 {ô}{{\^o}}1 {û}{{\^u}}1 {î}{{\^i}}1
}
\lstset{style=py}

\author{Stefania Marcassa}
\institute{CY Cergy Paris Université --- L3 Économie}
\date{}


\title{Environnement et anatomie des données}
\subtitle{Analyse des données --- Séance 1}

\begin{document}

\maketitle

\begin{frame}{Où nous allons aujourd'hui}
\begin{enumerate}
  \item Le notebook : comment il fonctionne, et comment il trompe.
  \item Les quelques objets Python dont nous aurons besoin tout le semestre.
  \item Une première application économique : un panier de biens.
  \item L'anatomie d'un tableau de données --- la partie qui compte vraiment.
\end{enumerate}

\vspace{1.5em}
À la fin de la séance, vous devez pouvoir répondre à une question : \alert{de quoi une ligne de ce fichier est-elle la description ?}
\end{frame}

%=========================================================================
\section{Le notebook}
%=========================================================================

\begin{frame}{Google Colab, et rien d'autre}
Tout le cours se fait dans le navigateur.

\vspace{1em}
\begin{itemize}
  \item Aucune installation, aucune version de Python à gérer
  \item Aucun problème de chemin d'accès
  \item Un compte Google suffit
  \item Le code s'exécute sur une machine distante, pas sur la vôtre
\end{itemize}

\vspace{1.2em}
Le travail en local (Anaconda, VS Code) est possible mais \alert{non assisté} : les TD ne sont pas consacrés au dépannage d'installations.
\end{frame}

\begin{frame}{Premier notebook --- maintenant}
\begin{enumerate}
  \item Connectez-vous à votre compte Google.
  \item Allez sur \texttt{colab.research.google.com}
  \item Cliquez sur \textbf{Nouveau notebook}.
\end{enumerate}

\vspace{1.5em}
Un notebook est une suite de \textbf{cellules} de deux types :

\vspace{0.6em}
\begin{itemize}
  \item les cellules de \textbf{code}, exécutées avec \texttt{Maj + Entrée} ;
  \item les cellules de \textbf{texte}, pour dire ce que l'on fait et pourquoi.
\end{itemize}

\vspace{1.2em}
Un notebook sans cellules de texte est illisible --- y compris pour vous, dans trois semaines.
\end{frame}

\begin{frame}[fragile]{Une première cellule}
\begin{lstlisting}
# Le taux de croissance entre deux annees
pib_2023 = 2822
pib_2024 = 2891

croissance = (pib_2024 - pib_2023) / pib_2023
print(croissance)
\end{lstlisting}

\vspace{0.5em}
\begin{lstlisting}
0.024450035435861798
\end{lstlisting}

\vspace{1em}
Trois choses : le \texttt{\#} introduit un commentaire, le \texttt{=} affecte une valeur, et rien ne s'affiche sans \texttt{print()}.
\end{frame}

\begin{frame}[fragile]{Le piège de l'ordre d'exécution}
Les cellules gardent leur effet même après modification. L'ordre dans lequel vous les avez \emph{exécutées} n'est pas l'ordre dans lequel elles sont \emph{écrites}.

\vspace{1em}
\begin{lstlisting}
tva = 0.20              # cellule executee en premier
prix_ttc = 100 * (1 + tva)

tva = 0.055             # cellule ajoutee plus tard, plus haut
\end{lstlisting}

\vspace{0.8em}
\texttt{prix\_ttc} vaut toujours 120. Rien ne le signale.

\vspace{1em}
\begin{alertblock}{Le réflexe à prendre}
Avant de conclure quoi que ce soit : \textbf{Exécution $\rightarrow$ Tout exécuter}. Si le notebook ne passe pas de haut en bas, il est faux.
\end{alertblock}
\end{frame}

\begin{frame}{Un piège qui vous coûtera dix minutes}
Lors du premier import de fichier, Colab peut renvoyer une erreur si les \textbf{cookies tiers} sont bloqués dans votre navigateur.

\vspace{1.2em}
La correction :
\begin{enumerate}
  \item Ouvrir les paramètres du navigateur
  \item Autoriser \texttt{https://[*.]googleusercontent.com:443}
\end{enumerate}

\vspace{1.2em}
Notez-le maintenant. Vous en aurez besoin à la séance 2.
\end{frame}

%=========================================================================
\section{Les objets dont nous aurons besoin}
%=========================================================================

\begin{frame}[fragile]{Quatre types de base}
\begin{lstlisting}
inflation = 2.4          # float  : nombre a virgule
annee     = 2024         # int    : nombre entier
pays      = "France"     # str    : texte
zone_euro = True         # bool   : vrai / faux

type(inflation)          # <class 'float'>
\end{lstlisting}

\vspace{1.5em}
La fonction \texttt{type()} paraît inutile aujourd'hui. Elle sera votre premier réflexe de diagnostic à partir de la séance 2.
\end{frame}

\begin{frame}{Pourquoi le type n'est pas un détail}
\begin{block}{Le code commune}
Le code INSEE de Bourg-en-Bresse est \texttt{01053}. Lu comme un nombre, il devient 1053 --- et ne correspond plus à aucune commune lors d'une fusion de fichiers.
\end{block}

\vspace{0.8em}
\begin{block}{L'année}
Une année stockée comme texte se trie dans l'ordre alphabétique. \texttt{"1999"} vient après \texttt{"10000"}, et les comparaisons \texttt{>} ne veulent plus rien dire.
\end{block}

\vspace{1.2em}
La moitié des bugs de ce semestre viendra d'un type mal interprété à la lecture d'un fichier. Aucun ne produira de message d'erreur.
\end{frame}

\begin{frame}[fragile]{Listes : une suite ordonnée}
\begin{lstlisting}
prix = [1.20, 0.95, 2.40, 1.75]

prix[0]        # 1.2   -- on compte a partir de zero
prix[-1]       # 1.75  -- le dernier
prix[0:2]      # [1.2, 0.95]  -- borne de droite exclue
len(prix)      # 4
sum(prix)      # 6.3
\end{lstlisting}

\vspace{1em}
Deux conventions à intégrer une fois pour toutes : l'indexation commence à \alert{zéro}, et la borne de droite d'une tranche est \alert{exclue}.

\vspace{0.6em}
Se tromper sur l'une ou l'autre ne produit pas d'erreur. Cela produit un résultat faux.
\end{frame}

\begin{frame}[fragile]{Dictionnaires : des paires clé--valeur}
\begin{lstlisting}
prix = {"pain": 1.20, "lait": 0.95, "cafe": 2.40}

prix["pain"]            # 1.2
prix.keys()             # les biens
prix.values()           # les prix
"beurre" in prix        # False
\end{lstlisting}

\vspace{1.2em}
Là où une liste retient un \emph{ordre}, un dictionnaire retient une \emph{correspondance}.

\vspace{0.8em}
C'est la structure d'une ligne de tableau : chaque variable porte un nom, et à chaque nom correspond une valeur.
\end{frame}

\begin{frame}[fragile]{Lire un message d'erreur}
\begin{lstlisting}
prix["beurre"]

KeyError: 'beurre'
\end{lstlisting}

\vspace{1em}
Un message d'erreur Python se lit \alert{par la fin} : la dernière ligne dit ce qui ne va pas, celles du dessus disent où.

\vspace{1em}
\begin{itemize}
  \item \texttt{KeyError} --- cette clé n'existe pas
  \item \texttt{NameError} --- variable jamais définie (cellule non exécutée ?)
  \item \texttt{TypeError} --- opération impossible entre ces types
  \item \texttt{FileNotFoundError} --- le fichier n'est pas là où vous croyez
\end{itemize}

\vspace{0.8em}
Une erreur qui s'affiche est une bonne nouvelle. Ce sont les calculs silencieusement faux qui coûtent cher.
\end{frame}

%=========================================================================
\section{Une première application}
%=========================================================================

\begin{frame}[fragile]{Un panier de biens}
Deux dictionnaires suffisent à représenter le problème économique.

\vspace{1em}
\begin{lstlisting}
prix_2023 = {"pain": 1.20, "lait": 0.95, "cafe": 2.40}
prix_2024 = {"pain": 1.35, "lait": 1.02, "cafe": 2.28}

quantites_2023 = {"pain": 120, "lait": 90, "cafe": 24}
\end{lstlisting}

\vspace{1.2em}
Notez que le prix du café \emph{baisse}. Un indice de prix n'est pas la moyenne des hausses : c'est le coût d'un panier donné, comparé dans le temps.
\end{frame}

\begin{frame}[fragile]{Le coût du panier}
\begin{lstlisting}
depense_2023 = 0
for bien in quantites_2023:
    depense_2023 += prix_2023[bien] * quantites_2023[bien]

print(depense_2023)     # 287.1
\end{lstlisting}

\vspace{1.2em}
Une boucle \texttt{for} parcourt les clés du dictionnaire ; le \texttt{+=} accumule. L'indentation n'est pas décorative : elle délimite le corps de la boucle.

\vspace{1em}
Refaites le calcul avec \texttt{prix\_2024} et les \emph{mêmes} quantités. Le rapport des deux dépenses est un indice de Laspeyres.
\end{frame}

\begin{frame}{La question qui occupera le TD}
Nous venons de fixer les quantités à celles de 2023. Rien ne nous y obligeait.

\vspace{1.5em}
\begin{columns}[T]
\begin{column}{0.46\textwidth}
\textbf{Laspeyres}

\vspace{0.5em}
Quantités de l'année de base.

\vspace{0.5em}
\small\color{CYgris}
Combien coûterait aujourd'hui le panier d'hier ?
\end{column}
\begin{column}{0.46\textwidth}
\textbf{Paasche}

\vspace{0.5em}
Quantités de l'année courante.

\vspace{0.5em}
\small\color{CYgris}
Combien coûtait hier le panier d'aujourd'hui ?
\end{column}
\end{columns}

\vspace{1.8em}
Les deux mesurent « l'inflation ». Ils ne donnent pas le même chiffre, et l'écart n'est pas une erreur : les ménages substituent.
\end{frame}

%=========================================================================
\section{Anatomie d'un tableau de données}
%=========================================================================

\begin{frame}{Trois mots à ne jamais confondre}
\begin{center}
\begin{tabular}{lll}
\toprule
\textbf{id\_menage} & \textbf{revenu} & \textbf{region} \\
\midrule
1001 & 32\,400 & Île-de-France \\
1002 & 28\,900 & Occitanie \\
1003 & 41\,200 & Bretagne \\
\bottomrule
\end{tabular}
\end{center}

\vspace{1.5em}
\begin{itemize}
  \item Une \textbf{observation} : une ligne. Ici, un ménage.
  \item Une \textbf{variable} : une colonne. Ici, le revenu.
  \item L'\textbf{unité statistique} : ce que représente une ligne. Ici, le ménage --- et non l'individu.
\end{itemize}
\end{frame}

\begin{frame}{Pourquoi l'unité statistique est le premier piège}
La même question économique change de réponse selon l'unité retenue.

\vspace{1.2em}
\begin{block}{Un exemple}
« Quel est le revenu moyen ? »

\vspace{0.5em}
Par \textbf{ménage} : un chiffre. Par \textbf{individu} : un autre. Par \textbf{unité de consommation} : un troisième.

\vspace{0.5em}
Aucun n'est faux. Ils ne répondent simplement pas à la même question.
\end{block}

\vspace{1.2em}
Avant toute analyse : \alert{de quoi une ligne est-elle la description ?}
\end{frame}

\begin{frame}{Trois structures de données}
\begin{columns}[T]
\begin{column}{0.31\textwidth}
\textbf{Coupe transversale}

\vspace{0.5em}
Plusieurs unités, une date.

\vspace{0.5em}
\small\color{CYgris}
Les salaires de 40\,000 individus en 2023.
\end{column}
\begin{column}{0.31\textwidth}
\textbf{Série temporelle}

\vspace{0.5em}
Une unité, plusieurs dates.

\vspace{0.5em}
\small\color{CYgris}
Le taux de chômage français, 1975--2025.
\end{column}
\begin{column}{0.31\textwidth}
\textbf{Panel}

\vspace{0.5em}
Plusieurs unités, plusieurs dates.

\vspace{0.5em}
\small\color{CYgris}
Le PIB de 27 pays, 2000--2024.
\end{column}
\end{columns}

\vspace{2em}
La structure détermine ce que l'on peut estimer --- et ce que l'on ne peut pas. Nous construirons un panel à la séance~4.
\end{frame}

\begin{frame}{Exercice éclair}
Pour chacun de ces jeux de données : quelle est l'unité statistique, et quelle est la structure ?

\vspace{1.5em}
\begin{enumerate}
  \item Le prix de clôture du CAC 40, chaque jour depuis 1990.
  \item Les réponses à l'Enquête Emploi du deuxième trimestre 2024.
  \item Les transactions immobilières enregistrées en France de 2014 à 2024.
  \item Le taux d'emploi des 27 pays de l'UE, chaque année depuis 2000.
\end{enumerate}

\vspace{1.5em}
\small\color{CYgris}
Le cas 3 est moins évident qu'il n'y paraît : discutons-en.
\end{frame}

%=========================================================================
\section{Pour la suite}
%=========================================================================

\begin{frame}{Avant la séance 2}
\begin{enumerate}
  \item Reprenez le calcul du panier et obtenez l'indice de Laspeyres pour 2024.
  \item Vérifiez que votre notebook s'exécute entièrement de haut en bas.
  \item Autorisez \texttt{googleusercontent.com} dans votre navigateur.
  \item Récupérez le notebook de la séance 2 sur le site du cours.
\end{enumerate}

\vspace{2em}
\begin{center}
\large \texttt{stefaniamarcassa.github.io/analyse\_des\_donnees}
\end{center}
\end{frame}

\begin{frame}[standout]
Des questions ?
\end{frame}

\end{document}
